% The soil-zone interface iteration in detail, negative recharge and the MODFLOW side
\subsection{The interface iteration in detail}\label{sec:szinterface}

\paragraph{Schedule.} After each of the first $n_{\mathrm{early}}=\min(5,\max(1,n_{\max}/5))$ outer iterations the interface
is updated whether or not MODFLOW has converged ($n_{\max}$ the outer-iteration limit, \cmd{:SolverMaxOuterIterations} or the
existing model's \cmd{OUTER\_MAXIMUM}); later only after outer iterations on which MODFLOW has converged, at most
\cmd{:SoilZoneIteration} \emph{max} of them. The last tenth of the limit (at least two outer iterations) is left to MODFLOW
on the boundary it has: the update after iteration $n_{\max}-\max(2,n_{\max}/10)$ and the last one the budget allows are
\emph{supply-limited} (below), and once the interface is frozen nothing changes. A step without any update is reported
unconverged. \cmd{:SoilZoneIteration 0} is explicit coupling (the predictor only), and so is an outer-iteration limit of 1,
which leaves no update; both report the interface converged. With a limit of 2 or 3 the step has a single update, whose
flow test compares with the predictor's total exchange.

\paragraph{The rejected recharge as a head-dependent term.} With the connection head $h_c$ of the update MODFLOW receives
the link's recharge $R_l\,(1-F(h_c))$ as a rate and, combined with the exchange's line, the linearization of the rejection
about $h_c$,
\begin{equation}
q_{\mathrm{rej}}(h)=s\,(h_c-h),\qquad s=\max\Big(R_lF'(h_c),\ \frac{R_l\big(F(h^*)-F(h_c)\big)}{h^*-h_c}\Big)\ge0,
\end{equation}
$h^*$ the head MODFLOW returned last (a GHB of head $h_c$ and conductance $s$, merged with the exchange's GHB). MODFLOW's own
iteration then sees that more recharge is rejected as its head rises (over a top layer with $K_v<0.01$~m/d, confined above
its top, the recharge is otherwise a valve of very high gain; with the term the interface of test S1 converges in 4 to 6 outer
iterations per step). The term is booked with the exchange, so mass is exact. On a link
whose unsaturated column carries the step the saturated share is rejected with the fraction of the start of the step $F_0$;
the water the column gives back to the soil zone on ground that saturated since then (rule (i)) follows the connection head
and is in the term, $R_q\min(\max(F(h)-F_0,0),1-F_0)$ with $R_q$ the column's outputs over $1-F_0$ (everything when $F_0=1$;
in explicit coupling only that).

\paragraph{The kink drain.} Where water rises into a soil zone without free water, nothing rises below the level where the
contact ends and the column has nothing to give: the exchange has a kink there, across which a line overshoots (it reverses
into flow out of a column that holds none) and the iteration cycles (test S2). MODFLOW then gets the exchange as a drain
(\cmd{SZD\_RAVEN}) ending at the line's zero, at least the cell bottom $+1$~mm (no DEM sample base lies below that level:
bases are clipped there), steepened at that level to carry the local solve's exchange at the head of linearization (at most
the larger of $C_g$ and the link conductance), its onset smoothed over at least 1~mm. A column that holds free water keeps
the line (it drains below the line's zero).

\paragraph{Convergence.} An update is converged when the last change of every column's soil level is below $\epsilon_h$,
the total exchange changed by no more than $\max(\mathrm{tol}_Q\,G,\ \sum F C\max(\epsilon_h,\mathrm{DVCLOSE}))$ and, per link
in head units (over $K=\max(C,S_ya/\Delta t)$): the rejection at MODFLOW's head and at the connection MODFLOW solved with
differ by less than $\epsilon_h$; the boundary MODFLOW solved with (its HCOF and RHS as booked) differs at the current head
from the local solve's exchange plus the change of the rejection by less than $\epsilon_h$, with only two exemptions --- (i) a
line (not a drain) whose local solves' secant through a real head change (more than $\epsilon_h$, the connection held) has the
wrong sign: its exchange part up to $B=\max(10\,\%\ |Q_c|,\ \epsilon_hK)$; (ii) a kink drain within 1~mm of its level: the
part MODFLOW's smoothing of the drain's onset takes; and the connection MODFLOW solved with gives the converged exchange at
the current head (lag below $\epsilon_hC$ beyond $B$). Then two more tests: the boundary MODFLOW solved with takes from no
column more than it holds at the current heads (the \emph{supply test}), and no link's connection head lies further than
$\max(d_s,10\epsilon_h)$ from MODFLOW's head where the saturated fractions differ by more than $10^{-6}$ (the
\emph{connection test}; specified heads and HDRY skipped). A link failing the connection test has its regula falsi bracket
and its last residual dropped, so that its next update is the plain one ($h_c=h$; counted in \cmd{GWModelSummary.txt},
``connection-test drops''). On the first ten unconverged steps \cmd{Raven\_errors.txt} names the step by its end, the last
change of the soil level and its link (HRU and top active cell), the change of the total exchange and its tolerance --- or,
when both were within the tolerances, the test that failed: the shortfall [m$^3$] and the link of the largest, or the
connection head's distance from MODFLOW's head, its link and the saturated fractions at both heads. No update is made from an
iterate MODFLOW has not converged on that holds a head that is not finite, or in which a link's top cell lies far outside the
model's elevations (the divergence margin below) and away from that link's connection head (more than the margin from it,
and on the other side of the model or more than the margin further outside).

\paragraph{Supply limits.} In explicit coupling and on the supply-limited update the boundary MODFLOW keeps may not take more
from a column than its free water, down to the lowest head the cell can reach (the bottom of the column's lowest active
cell): first the rejection's term ($s(h_c-z_{\mathrm{low}})$ at most the head-dependent rejection and at most the free water
the line leaves), then the line with what the term leaves. \cmd{GWModelSummary.txt} reports the number of such link updates
and the largest take beyond the free water (0 up to round-off).

\paragraph{Negative recharge.}\label{sec:negrecharge} A process that takes water out of \cmd{GROUNDWATER} (e.g. \cmd{PERC\_GR4JEXCH SOIL[1]
GROUNDWATER} with a positive \cmd{GR4J\_X2}) gives a negative recharge. On an HRU with a soil zone it is never split by $F$: all
of it is taken from MODFLOW (through the column on a column link), at most what the top cell can still give in the step
(half its water above its bottom under the link, less what the columns, the capillary rise and the repayment of a lagged
debt took before); the rest is cut and taken back out of the soil zone as a negative rejected recharge (run total ``negative
recharge cut''). What the soil zone does not hold of that cut (a process after the import moved it on before the coupled
exchange) is taken from the HRU's other columns in proportion to their water, and only what the whole HRU cannot supply is
booked as ``negative recharge the soil zone did not hold'' (its own \cmd{GWBudget.csv} column, warning and run total). On HRUs
without a soil zone and without \cmd{:SoilZone PHYSICAL} it is plain recharge, not limited (a warning).

\paragraph{Drained columns.}\label{sec:drained} With \cmd{NEWTON}, before each solve a column whose active cells are all convertible
(\cmd{ICELLTYPE}$\neq0$: a generated model's layer types, an existing model's from MODFLOW's memory) and whose lowest cell is
at or below its bottom is set 1~mm above that bottom in MODFLOW's iterate, once per step, once more on a later solve of the
step that leaves the lowest cell more than 1~m below its bottom, and after a solve with heads that are not finite; never in
the last two outer iterations (except for heads that are not finite). A column with an active confined cell is never drained.

\paragraph{Heads far outside the model.} A step whose heads are not finite stops the run at once. A step MODFLOW did not
converge on that holds a head beyond the model's highest top or lowest bottom by more than ten times its relief (at least
1~km) is carried on only if every such head that moved in the step (relative to the heads of the last converged step, by
position) moved by at most that margin in the last outer iteration and MODFLOW's residual is at most 1\,\% of the gross flow
through the aquifer (every entry's inflow and outflow, MODFLOW's TOTAL IN plus TOTAL OUT; Raven's percentage divides by
IN$+$OUT, half MODFLOW's PERCENT DISCREPANCY); such a step is warned with its residual. Otherwise the run stops, naming the
head, its cell, its last move, the residual and the usual causes. An unconverged step without a far head is carried on
whatever its residual, as MODFLOW's own \cmd{CONTINUE} option would: the residual is MODFLOW's balance error and a negative one
is withheld from the discharge returned to Raven. A far head on a converged step is warned about once per run (per ensemble
member).

\paragraph{Existing models.} \cmd{:SoilZone PHYSICAL} needs an existing model with \cmd{NEWTON} (start-up stop naming the name
file; MODFLOW then needs \cmd{LINEAR\_ACCELERATION BICGSTAB}). A model without a storage package is steady state: storage 0,
$S_y$ taken from \cmd{SPEC\_YIELD} of \cmd{:SoilZoneAquiferClass} when that class gives it, otherwise 0.2 (a warning naming
the value either way) for the soil-zone checks, the release limits and the columns; with one, \cmd{SY} is taken as
MODFLOW holds it, also where it is 0 (common with \cmd{ICONVERT} 0, all of the model included), floored at 0.01 for the
release limits and the columns, with a warning counting the cells below 0.01: such a model's outputs are unchanged. The warning
about an outer-iteration limit below 20 is written once MODFLOW's limit is known (\cmd{OUTER\_MAXIMUM} for an existing model).
A restart writes the heads with 17 significant digits and writes 0 as the lagged exchange of links whose cell was a specified
head in the last step. Without \cmd{NEWTON} (\cmd{:SoilZone CONCEPTUAL}) the recharge MODFLOW does not apply to a column whose
cells are all dry or whose highest active cell is a specified head is booked, signed, as ``recharge not applied by MODFLOW''
(the water such cells held when they dried is in neither balance; confined cells never dry).
